\chapter{Progettazione e codifica}
\label{chap:progettazione-codifica}

\textit{In questo capitolo vengono descritte le scelte tecnologiche, architetturali e implementative adottate durante lo sviluppo del sistema. Viene prima fornita una panoramica delle tecnologie e degli strumenti utilizzati, seguita dall'illustrazione dell'architettura generale, dei design pattern applicati e dei principali componenti realizzati.}

\section{Tecnologie e strumenti}
\label{sec:tecnologie-strumenti}

Il sistema è realizzato interamente in C++ ed è articolato in tre target applicativi che condividono un unico nucleo di pianificazione: un'applicazione desktop Qt usata come strumento di debug e testing durante lo sviluppo, un server HTTP Drogon come interfaccia REST di produzione verso l'ERP e un server gRPC opzionale per delegare la sola fase di calcolo a una macchina dedicata. La scelta delle tecnologie è stata orientata da requisiti di prestazione, integrazione con l'ecosistema ERP esistente e supporto alla programmazione a vincoli.

\subsection{Linguaggi e strumenti di sviluppo}
\label{subsec:linguaggi}

\subsubsection{C++20}

C++20 è lo standard ISO del linguaggio C++ adottato nel progetto (requisito RVN-1). La scelta di C++ come linguaggio principale è dettata anche dalla necessità di integrarsi direttamente con la libreria Google OR-Tools, la cui API nativa è in C++.

\subsubsection{Qt6}

Qt è un framework multipiattaforma per lo sviluppo di applicazioni desktop, mantenuto da Qt Group. Qt6 è la versione adottata per il target \texttt{erpai\_ui}, l'applicazione desktop sviluppata principalmente come strumento di debug e testing durante le sessioni di sviluppo. Grazie alla sua interfaccia grafica interattiva, permette di configurare e avviare pianificazioni, visualizzare il diagramma di Gantt, i KPI e i log del solver, e riprodurre scenari di test senza ricorrere alla pipeline HTTP di produzione. Tra le caratteristiche fondamentali di Qt6 per il progetto:

\begin{itemize}
    \item \textbf{Signals \& Slots}: meccanismo di comunicazione tra oggetti che consente di connettere eventi (\textit{signals}) emessi da un oggetto a metodi (\textit{slots}) di un altro, anche tra thread diversi, grazie alle \textit{queued connections} che garantiscono la thread-safety degli aggiornamenti alla UI.
    \item \textbf{QThread e pattern worker}: permettono di eseguire operazioni lunghe in background senza bloccare il thread principale della UI, pattern fondamentale per integrare il solver CP-SAT nell'applicazione desktop.
    \item \textbf{Widget predefiniti}: \texttt{QTableWidget}, \texttt{QPlainTextEdit}, \texttt{QSplitter} e altri componenti grafici standard usati per costruire le viste di configurazione, monitoraggio e risultati.
    \item \textbf{QPainter}: classe per il disegno vettoriale 2D, impiegata per il rendering personalizzato del diagramma di Gantt nella classe \texttt{GanttWidget}.
\end{itemize}

\subsubsection{CMake}

CMake è un sistema di build open source e multipiattaforma che genera file di build nativi (Makefile, Ninja, progetti Visual Studio) a partire da una descrizione dichiarativa in file \texttt{CMakeLists.txt}. Nel progetto CMake gestisce la compilazione dei tre target principali (\texttt{erpai\_ui}, \texttt{erpai\_grpc\_solver}, \texttt{erpai\_http}) e le dipendenze verso OR-Tools, MariaDB Connector/C++, Qt6, Drogon e le librerie gRPC/Protobuf.

\subsection{Librerie}
\label{subsec:librerie}

\subsubsection{Google OR-Tools}

Google OR-Tools è una libreria open source sviluppata da Google per la risoluzione di problemi di ottimizzazione combinatoria. Il componente chiave utilizzato nel progetto è il solver \textbf{CP-SAT} (\textit{Constraint Programming - Satisfiability}), un motore di programmazione a vincoli basato su tecniche SAT moderne. CP-SAT è particolarmente adatto alla schedulazione di produzione perché supporta nativamente:

\begin{itemize}
    \item \textbf{Variabili di intervallo} (\texttt{IntervalVar}): rappresentano i task come intervalli temporali con variabili di inizio, fine e durata.
    \item \textbf{Vincoli di non sovrapposizione} (\texttt{AddNoOverlap}): garantiscono che su ogni macchina non vengano eseguiti due task contemporaneamente.
    \item \textbf{Vincoli di capacità cumulativa} (\texttt{AddCumulative}): modellano la capacità dei gruppi di macchine, in cui più macchine dello stesso gruppo possono lavorare in parallelo.
    \item \textbf{Funzioni obiettivo multi-criterio}: il solver minimizza la lateness totale rispetto alle date di consegna, usando il makespan come criterio secondario di spareggio.
\end{itemize}

Il solver opera su dominio intero, coerente con la rappresentazione dei tempi in minuti adottata nel progetto. La scelta di OR-Tools è dettata dal requisito RVN-3.

\subsubsection{MariaDB Connector/C++}

MariaDB Connector/C++ è il driver ufficiale per la connessione a database MariaDB/MySQL da applicazioni C++. Viene impiegato per caricare dal database ERP i dati di input alla pianificazione: ordini di produzione (\texttt{mrp\_steps\_plans}), fasi di lavorazione (\texttt{steps}), macchine, gruppi macchina e relative associazioni. Le strutture dati caricate (es. \texttt{MrpStepsPlans}, \texttt{Steps}, \texttt{Machines}) sono modellate come classi C++ corrispondenti alle tabelle del gestionale aziendale, gestite nel namespace \texttt{db}.

\subsubsection{gRPC e Protocol Buffers}

gRPC è un framework open source per chiamate a procedure remote ad alte prestazioni, sviluppato da Google. Utilizza Protocol Buffers come formato di serializzazione binaria dei messaggi.

Nel progetto, gRPC non è un'interfaccia generica: è l'opzione di deployment che consente di delegare \textbf{esclusivamente la fase di calcolo CP-SAT} a un server dedicato. Il target \texttt{erpai\_grpc\_solver} espone il servizio \texttt{PlanningService} con il metodo \texttt{RunPlanning}, che riceve il modello di scheduling già costruito dal chiamante e ne esegue la risoluzione, restituendo in streaming gli aggiornamenti di stato (KPI, log, soluzione incumbent).

In questo scenario, il modulo chiamante --- che sia il server HTTP, il daemon o l'applicazione Qt --- continua a eseguire internamente tutte le fasi a monte: caricamento dei dati da MariaDB, costruzione del modello CP-SAT tramite \texttt{ModelBuilder} e orchestrazione del flusso tramite \texttt{PlanningRunner}. La delega al server gRPC riguarda unicamente la parte computazionalmente più onerosa, permettendo di isolarla su una macchina ad alte prestazioni senza modificare la logica applicativa del chiamante. La classe \texttt{RemoteSolverClient} incapsula questa comunicazione in modo trasparente.

\subsubsection{Drogon}

Drogon è un framework C++ per la realizzazione di server HTTP e WebSocket ad alte prestazioni. Nel progetto viene impiegato per il target \texttt{erpai\_http}, l'interfaccia REST di produzione attraverso cui l'applicazione ERP in PHP comunica con il solver. Le route esposte permettono di avviare una pianificazione, leggere il piano risultante e consultare lo storico delle esecuzioni. Grazie alla sua architettura asincrona e non bloccante basata su coroutine, Drogon gestisce richieste concorrenti senza penalizzare il tempo di risposta.

\subsection{Containerizzazione}
\label{subsec:docker}

\subsubsection{Docker}

Docker è una piattaforma open source per la containerizzazione di applicazioni. I target server del progetto (gRPC e HTTP) vengono distribuiti come immagini Docker tramite i file \texttt{Dockerfile.grpc} e \texttt{Dockerfile.http}, accompagnati dai rispettivi file \texttt{docker-compose.grpc.yml} e \texttt{docker-compose.http.yml}. Questo approccio garantisce la riproducibilità dell'ambiente di esecuzione e semplifica il deploy del componente solver in scenari di integrazione con sistemi aziendali esterni.

\subsection{Persistenza locale}
\label{subsec:sqlite}

\subsubsection{SQLite}

SQLite è una libreria C che implementa un motore di database \gls{sql} embedded, senza necessità di un server separato. Nel progetto viene impiegata per la persistenza dello storico locale delle pianificazioni (requisito RVN-5). La classe \texttt{LocalRunStore} gestisce un file \texttt{.db} nella directory dell'eseguibile, creando e mantenendo le tabelle \texttt{runs}, \texttt{run\_logs} e \texttt{run\_tasks} attraverso le API C di SQLite. La scelta di SQLite garantisce leggerezza, portabilità e assenza di dipendenze esterne lato client.

\section{Architettura del sistema}
\label{sec:architettura}

Il sistema è progettato attorno a un \textbf{nucleo condiviso} (\texttt{core}) separato da qualsiasi dipendenza verso UI, database o protocollo di comunicazione. Questo nucleo espone \texttt{PlanningRunner} come orchestratore del flusso di pianificazione e \texttt{IPlanningObserver} come punto di estensione per la distribuzione degli eventi.

Tre target applicativi condividono lo stesso nucleo:

\begin{itemize}
    \item \textbf{erpai\_ui}: applicazione desktop Qt, sviluppata come strumento di debug e testing durante lo sviluppo; fornisce una \gls{gui} interattiva per avviare pianificazioni e analizzarne i risultati in tempo reale.
    \item \textbf{erpai\_http}: server HTTP basato su Drogon, interfaccia REST di produzione per l'integrazione con l'applicazione ERP in PHP.
    \item \textbf{erpai\_grpc\_solver}: server gRPC opzionale che, se attivo, riceve il modello CP-SAT già costruito dal modulo chiamante ed esegue la sola fase di calcolo su una macchina dedicata.
\end{itemize}

Questa separazione garantisce che la logica di pianificazione possa essere testata, eseguita e integrata indipendentemente dal canale di accesso, rispettando il principio di separazione delle responsabilità. In entrambi i casi --- che il transport layer sia HTTP o gRPC --- il caricamento dei dati, la costruzione del modello e l'orchestrazione del flusso avvengono nel modulo principale; la variabile è solo come e dove viene eseguita la fase di risoluzione CP-SAT.

\subsection{Pipeline di pianificazione}
\label{subsec:pipeline}

Il flusso di esecuzione di una pianificazione segue le fasi descritte di seguito.

\begin{enumerate}
    \item \textbf{Caricamento dati} (\texttt{MrpDataLoader}): carica da MariaDB gli ordini di produzione, le fasi di lavorazione, le macchine, i calendari lavorativi per macchina e le impostazioni \gls{mrp}.
    \item \textbf{Costruzione task} (\texttt{PlanningInput}): trasforma i record di database in strutture \texttt{TaskData} con durata calcolata (\texttt{setup\_seconds + qta $\times$ unit\_seconds}) e relazioni di precedenza tra fasi dello stesso job, derivate dal \texttt{phase\_number}.
    \item \textbf{Costruzione modello CP-SAT} (\texttt{ModelBuilder}): crea per ogni task le variabili CP-SAT (\texttt{IntVar start}, \texttt{IntVar end}, \texttt{IntervalVar interval}), aggiunge i vincoli di non sovrapposizione per macchina, i vincoli di capacità cumulativa per gruppo macchina, le precedenze tra fasi e i vincoli di calendario lavorativo giornaliero.
    \item \textbf{Risoluzione} (\texttt{TwoPhaseSolver}): esegue la pipeline a due fasi descritta nella sezione~\ref{subsec:twophase}.
    \item \textbf{Estrazione risultati} (\texttt{ResultExtractor}): converte la soluzione CP-SAT in strutture \texttt{ScheduledTaskInfo} con inizio e fine in data/ora assoluta.
    \item \textbf{Notifica osservatori} (\texttt{IPlanningObserver}): distribuisce KPI, log e task schedulati a tutti gli osservatori registrati.
\end{enumerate}

\subsection{Solver a due fasi}
\label{subsec:twophase}

Il \texttt{TwoPhaseSolver} implementa una pipeline di risoluzione a due fasi per gestire la complessità del problema in presenza di calendari lavorativi e vincoli multipli.

\begin{itemize}
    \item \textbf{Fase 1 (Light)}: risolve il problema su un modello semplificato (\texttt{BuildMode::Phase1Light}), senza i vincoli di calendario più stringenti. Produce una soluzione iniziale in tempi brevi.
    \item \textbf{Hinting}: la soluzione della fase 1 viene estratta tramite \texttt{SolutionHintBuilder} e iniettata come suggerimento (\gls{warmstart}) nel modello completo, guidando il solver verso la regione dello spazio di ricerca più promettente.
    \item \textbf{Fase 2 (Full)}: ottimizza sul modello completo (\texttt{BuildMode::Phase2Full}) partendo dall'hint. Gestisce calendari giornalieri per macchina, vincoli di capacità cumulativa e funzione obiettivo multi-criterio (lateness totale dei job + makespan come criterio secondario).
\end{itemize}

Qualora la fase 2 non produca una soluzione migliore entro il timeout configurato, il sistema blocca il salvataggio a database. (\textit{fallback}).

\section{Design Pattern adottati}
\label{sec:design-pattern}

\subsection{Observer}
\label{subsec:observer}

Il pattern Observer definisce una dipendenza uno-a-molti tra oggetti, in modo che quando un oggetto cambia stato tutti i suoi osservatori vengano notificati automaticamente. Nel sistema sviluppato, il pattern è implementato tramite l'interfaccia astratta \texttt{IPlanningObserver} (namespace \texttt{core}).

L'interfaccia dichiara i metodi di notifica per tutti gli eventi rilevanti del ciclo di vita di una pianificazione: avvio del run (\texttt{onRunStarted}), cambio di fase (\texttt{onPhaseChanged}), messaggi di log (\texttt{onLogMessage}), aggiornamenti KPI (\texttt{onKpiUpdated}), task schedulati (\texttt{onTaskScheduled}), warning di validazione (\texttt{onValidationWarning}) e completamento del run (\texttt{onRunFinished}). Metodi aggiuntivi per la telemetria realtime del solver (\texttt{onSolverTelemetry}, \texttt{onSolverIncumbent}, \texttt{onTaskBatch}, ecc.) sono definiti con implementazione di default no-op, preservando la retro-compatibilità con le implementazioni esistenti.

\begin{listing}[H]
\begin{minted}{cpp}
class IPlanningObserver {
public:
    virtual ~IPlanningObserver() = default;

    virtual void onRunStarted(const std::string& plan_start,
                              int total_tasks) = 0;
    virtual void onPhaseChanged(const PhaseInfo& phase) = 0;
    virtual void onLogMessage(const std::string& message,
                              ValidationWarning::Level level
                                  = ValidationWarning::Level::Info) = 0;
    virtual void onKpiUpdated(const SolverKpi& kpi) = 0;
    virtual void onTaskScheduled(const ScheduledTaskInfo& task) = 0;
    virtual void onValidationWarning(const ValidationWarning& w) = 0;
    virtual void onRunFinished(const PlanningReport& report) = 0;

    // Metodi con implementazione di default no-op (retro-compatibili)
    virtual void onSolverTelemetry(const SolverTelemetrySnapshot&) {}
    virtual void onTaskBatch(const TaskBatchEvent&) {}
    virtual void onMachineTimeline(
        const std::vector<MachineTimelineRow>&) {}
};
\end{minted}
\caption{Interfaccia \texttt{IPlanningObserver} (semplificata).}
\label{listing:observer}
\end{listing}

Tre implementazioni concrete dell'interfaccia vengono utilizzate nel progetto:
\begin{itemize}
    \item \texttt{SqliteRunObserver}: persiste i dati del run nello storico SQLite locale.
    \item \texttt{HttpPlanningObserver}: invia callback HTTP al termine di ogni fase.
    \item L'implementazione Qt nel \texttt{PlannerWorker}: traduce le notifiche in segnali Qt, instradati alla UI tramite queued connections.
\end{itemize}

Il \texttt{CompositeObserver} consente di registrare più osservatori simultaneamente, effettuando il fan-out delle notifiche verso ciascuno.

\subsection{Strategy}
\label{subsec:strategy}

Il pattern Strategy definisce una famiglia di algoritmi intercambiabili, incapsulandoli in unità separate. Nel sistema, le diverse strategie di risoluzione CP-SAT sono implementate come funzioni distinte nello spazio dei nomi \texttt{planning}:

\begin{itemize}
    \item \texttt{solveBasicNoOverlap}: modello minimale con soli vincoli di non sovrapposizione per macchina.
    \item \texttt{solveWithDueDates}: aggiunge la penalizzazione del ritardo rispetto alle date di consegna.
    \item \texttt{solveWithGroupCapacityAndFixedOrders}: introduce i vincoli di capacità cumulativa per gruppi di macchine e gli ordini di produzione già confermati.
    \item \texttt{solveWithPhasePrecedences}: aggiunge i vincoli di precedenza tra le fasi dello stesso job.
    \item \texttt{solveTwoPhase}: pipeline a due fasi con hinting, descritta nella sezione~\ref{subsec:twophase}.
\end{itemize}

La scelta della strategia da applicare è controllata dall'enum \texttt{SolverPhases} e dalla struttura di configurazione \texttt{SolverRunConfig}, passata a \texttt{PlanningRunner} prima dell'avvio.

\subsection{Worker Thread}
\label{subsec:worker}

Nel contesto di Qt, il pattern Worker Thread separa l'esecuzione di operazioni potenzialmente lunghe dal thread principale della UI. La classe \texttt{PlannerWorker} (namespace \texttt{ui}) implementa questo pattern estendendo \texttt{QObject}.

\texttt{PlannerWorker} viene istanziato dalla \texttt{MainWindow}, spostato su un \texttt{QThread} dedicato tramite \texttt{moveToThread} e connesso ai widget della UI tramite signals/slots. I segnali emessi da \texttt{PlannerWorker} durante l'esecuzione del solver (es. \texttt{kpiUpdated}, \texttt{taskBatch}, \texttt{machineTimeline}, \texttt{logMessage}) vengono recapitati al thread della UI attraverso queued connections automatiche, garantendo la thread-safety degli aggiornamenti grafici senza alcun mutex esplicito nel codice UI.

\begin{listing}[H]
\begin{minted}{cpp}
class PlannerWorker : public QObject {
    Q_OBJECT
public slots:
    void startRunWithConfig(planning::SolverRunConfig cfg);
signals:
    void runStarted(const QString& plan_start, int total_tasks);
    void phaseChanged(const QString& name, int idx, const QString& desc);
    void logMessage(const QString& message, int level);
    void kpiUpdated(core::SolverKpi kpi);
    void taskBatch(core::TaskBatchEvent batch);
    void machineTimeline(std::vector<core::MachineTimelineRow> rows);
    void runFinished(core::PlanningReport report);
    void finished();
};
\end{minted}
\caption{Interfaccia pubblica di \texttt{PlannerWorker} (semplificata).}
\label{listing:plannerworker}
\end{listing}

\subsection{Repository}
\label{subsec:repository}

Il pattern Repository fornisce un'astrazione tra la logica applicativa e il layer di persistenza dei dati. L'interfaccia \texttt{RunStoreAccess} definisce il contratto per l'accesso allo storico locale. La classe \texttt{LocalRunStore} implementa tale contratto utilizzando SQLite come motore di persistenza.

\texttt{LocalRunStore} gestisce tre tabelle: \texttt{runs} (metadati e KPI di ogni pianificazione), \texttt{run\_logs} (messaggi di log associati al run) e \texttt{run\_tasks} (task schedulati con inizio, fine e macchina assegnata). Tutte le operazioni di scrittura e lettura sono protette da un \texttt{std::mutex}, garantendo la correttezza in caso di accesso concorrente da più thread.

\section{Struttura del progetto}
\label{sec:struttura-progetto}

Il codice sorgente è organizzato in moduli con responsabilità ben definite, riportati nella struttura di cartelle seguente.

\begin{listing}[H]
\begin{minted}{text}
src/
+-- apps/
|   +-- grpc_server/  # target server gRPC
|   +-- http/         # target server HTTP REST
|   +-- ui/           # target applicazione desktop Qt
|   |   +-- monitor/  # campionamento risorse di processo
|   |   +-- widgets/  # MainWindow, GanttWidget, SparklineChart
|   |   '-- worker/   # PlannerWorker (thread separato)
+-- core/
|   +-- observer/     # IPlanningObserver, CompositeObserver, PlanningEvents
|   '-- PlanningRunner.h/.cpp
+-- db/               # strutture dati ORM (tabelle ERP)
+-- local_db/         # SqliteRunObserver, LocalRunStore, RunStoreAccess
+-- planning/
|   +-- mrp/          # MrpDataLoader, BomRecursion, ForecastBatching
|   +-- planner/      # ModelBuilder, TwoPhaseSolver, SolverPhases
|   '-- reporting/    # ConsoleScheduleReport
+-- proto/            # solver_service.proto (contratto gRPC)
'-- utils/            # DateTime e utility generali
\end{minted}
\caption{Struttura delle cartelle del progetto.}
\label{listing:struttura}
\end{listing}

La separazione tra \texttt{core}, \texttt{planning} e \texttt{apps} riflette il principio di dipendenza unidirezionale: le applicazioni dipendono dal core, il core non dipende da nessuna applicazione. In questo modo è possibile aggiungere nuovi target (es. un client \gls{cli}) senza modificare la logica di pianificazione.

\section{Componenti principali}
\label{sec:componenti}

\subsection{PlanningRunner}
\label{subsec:planningrunner}

\texttt{PlanningRunner} è l'orchestratore centrale del flusso di pianificazione. Espone due metodi \texttt{run}: uno semplificato che accetta data di inizio e timeout, e uno avanzato che riceve una struttura \texttt{SolverRunConfig} con tutti i parametri configurabili. Prima dell'avvio, l'osservatore viene impostato tramite \texttt{setObserver}; il runner invoca i callback dell'observer ad ogni fase significativa del flusso.

\subsection{ModelBuilder}
\label{subsec:modelbuilder}

\texttt{ModelBuilder} costruisce il modello CP-SAT completo per la pianificazione. La funzione principale \texttt{buildGroupCapacityModel} crea le variabili di intervallo per ogni task o chunk di task, applica i vincoli di non sovrapposizione sulle macchine (\texttt{AddNoOverlap}), i vincoli di capacità cumulativa sui gruppi di macchine (\texttt{AddCumulative}), i vincoli di precedenza tra fasi dello stesso job e i vincoli del calendario lavorativo giornaliero per macchina. La funzione \texttt{applyStandardObjective} aggiunge la funzione obiettivo al modello: minimizzazione della lateness totale dei job se sono presenti date di consegna, altrimenti minimizzazione del makespan.

\subsection{TwoPhaseSolver}
\label{subsec:twophasesolver}

\texttt{TwoPhaseSolver} implementa la funzione \texttt{solveTwoPhase}, che coordina le due fasi di risoluzione descritte nella sezione~\ref{subsec:twophase}. Notifica l'observer ad ogni cambio di fase e al termine di ciascuna risoluzione, consentendo all'interfaccia grafica di mostrare lo stato di avanzamento in tempo reale. La struttura \texttt{SchedulerBenchResult} registra per ogni esecuzione lo stato del solver (\texttt{OPTIMAL}, \texttt{FEASIBLE}, \texttt{INFEASIBLE}), il valore dell'obiettivo, il bound migliore e il tempo impiegato; la funzione \texttt{selectBestBenchResult} determina il miglior risultato secondo una gerarchia di criteri: ottimalità, obiettivo minore, gap assoluto minore, tempo minore.

\subsection{LocalRunStore}
\label{subsec:localrunstore}

\texttt{LocalRunStore} gestisce la persistenza dello storico locale su file SQLite. L'inserimento di un run avviene in tre fasi: \texttt{insertRun} (status \textit{running}), \texttt{updateRunKpi} (aggiornamenti intermedi durante il run) e \texttt{finalizeRun} (chiusura con KPI finali e status \textit{finished} o \textit{failed}). I log vengono scritti incrementalmente con \texttt{insertLog}; i task schedulati vengono salvati in batch con \texttt{insertRunTasks} al termine del run. Tutte le operazioni sono thread-safe grazie a un \texttt{std::mutex} interno.

\subsection{PlannerWorker}
\label{subsec:plannerworker}

\texttt{PlannerWorker} è il bridge tra il solver CP-SAT e la UI Qt. Implementa l'interfaccia \texttt{IPlanningObserver} traducendo ogni callback del solver in un segnale Qt. Poiché il solver gira in un thread separato, le queued connections di Qt garantiscono che i segnali vengano processati nel thread della UI senza necessità di sincronizzazione esplicita. I tipi custom (es. \texttt{SolverKpi}, \texttt{PlanningReport}, \texttt{MachineTimelineRow}) sono registrati nel meta-object system di Qt tramite \texttt{Q\_DECLARE\_METATYPE}.

\subsection{GanttWidget}
\label{subsec:ganttwidget}

\texttt{GanttWidget} è un widget Qt personalizzato che visualizza il diagramma di Gantt del piano generato. Il rendering è realizzato con \texttt{QPainter}: sull'asse orizzontale è rappresentata la linea temporale del piano; sull'asse verticale sono elencate le macchine. Ogni task schedulato è disegnato come un blocco colorato sulla riga della macchina assegnata, con tooltip che mostrano il dettaglio dell'ordine. 

\newpage
